\section*{Related Work}

\subsection*{Archaic introgression detection: allele-frequency and density-based methods}

The detection of archaic hominin introgression in modern human genomes was first
made possible by comparative analysis of the initial Neanderthal and Denisovan
genome sequences~\citep{green2010draft, reich2010genetic}.
Early methods identified introgressed segments by searching for genomic regions
in non-African populations that showed excess similarity to the archaic genome,
quantified through D-statistics (ABBA-BABA tests) or derived allele sharing
frequency~\citep{patterson2012ancient, durand2011testing, green2010draft}.
While powerful for genome-wide proportion estimation, these population-level statistics
do not localize introgressed segments in individual genomes.

The first generation of segment-level methods took a density-based approach.
Vernot \& Akey~\citeyear{vernot2014resurrecting} used a conditional random field
that combined archaic allele density, linkage disequilibrium, and population
differentiation to map Neanderthal segments in European genomes; their method
identified $\sim$20~Gb of Neanderthal-derived sequence across the European
population.
Sankararaman et al.~\citeyear{sankararaman2014genomic} independently developed
a two-state HMM in which each window is classified as introgressed or not
based on the density of archaic-like alleles relative to a background model
calibrated from African genomes.
Both methods operationalize the same core statistical signal: the local excess of
archaic allele matches over a background expectation.
Neither explicitly models the spatial autocorrelation of archaic matches along
the haplotype.

IBDmix~\citep{chen2020identifying} reformulated the problem as identity-by-descent
detection between a query individual and the archaic genome, using a LOD score
at each site and a simple HMM to call IBD segments.
IBDmix was designed for sensitivity at the diploid-individual level and has
been widely used for large-scale population surveys (e.g., the HGDP--1000GP
combined dataset~\citep{browning2018ancestry}).
Its site-level LOD accumulation is equivalent to a Poisson approximation in the
number of matching sites, making it conceptually aligned with the Poisson density
principle our baseline captures.
Across this spectrum, the unifying signal remains the local excess of archaic
allele matches --- the density principle --- rather than the spatial structure
of haplotype matches that coalescent theory predicts to differ sharply between
introgressed and non-introgressed tracts.

\subsection*{HMM-based density methods: HMMix and DAIseg}

HMMix~\citep{skov2018detecting} represents the current state of the art in
HMM-based archaic introgression detection.
It uses a two-state HMM (introgressed vs.\ non-introgressed) with Poisson-distributed
emission probabilities for the count of archaic-diagnostic SNPs per window.
The transition matrix encodes the expected tract length through a geometrically
distributed segment length prior.
HMMix was used to produce the most comprehensive published archaic introgression
catalog in non-African populations, covering all major population groups.

DAIseg~\citep{planche2024daiseg} (Planche et al.\ 2024) extended HMMix to a three-state model (AMH/NEA/DEN),
enabling simultaneous detection of both archaic ancestry sources.
Like HMMix, DAIseg uses Poisson-distributed emissions at the site level.
The key structural difference from \AP{} is that both HMMix and DAIseg count
the number of archaic-diagnostic variants within a window and model this count
as Poisson-distributed, rather than modeling the per-site match probability
conditioned on a haplotype state.
This means neither method leverages the spatial match-length signal that
differentiates introgressed from non-introgressed haplotypes under the coalescent.
\AP{} is structurally closest to DAIseg in having three states and supporting
simultaneous NEA and DEN attribution, but differs fundamentally in its emission
model: where DAIseg counts variants per window, \AP{} computes a per-site
match probability using coalescent-theoretic parameters, explicitly encoding
the expected haplotype match length in the transition matrix.

\subsection*{Haplotype-based approaches: S'PRIME, ArchIE, and local ancestry inference}

Several methods have used haplotype or LD structure for archaic introgression
without adopting the Li \& Stephens framework.
S'PRIME~\citep{browning2018ancestry} identifies candidate introgressed segments
as regions with elevated frequency of derived alleles absent from African
genomes, combined with LD-based selection to maximize the number of candidate
introgressed SNPs in each putative segment.
While LD-aware, S'PRIME uses LD as a segment-extension heuristic rather than
as an explicit emission model.

ArchIE~\citep{durvasula2019statistical} uses a logistic regression classifier trained on
haplotype summary statistics (including match-length-related features) to classify
genomic windows as introgressed or not.
ArchIE demonstrated that haplotype features improve detection accuracy over
site-frequency features alone, directly motivating the present work.
However, its machine learning framework requires simulated training data, and
Peede et al.~\citeyear{peede2025review} showed that such simulation-trained
methods are sensitive to demographic misspecification when the training
demography departs from the true history.
\AP{} avoids this dependence by deriving its parameters from coalescent theory
directly, requiring only published admixture time estimates rather than
simulation-tuned weights.

The Li \& Stephens (LS) model~\citep{li2003modeling} is the theoretical backbone
of modern local ancestry inference (LAI).
In the LS framework, a query haplotype is modeled as an imperfect mosaic of
reference panel haplotypes; the HMM state at each site tracks which reference
haplotype is currently being ``copied'', with transitions governed by recombination
distance.
ChromoPainter~\citep{lawson2012inference} applies the LS model to population
structure and admixture inference, and MOSAIC~\citep{salter2019mosaic} extends
it to multi-way admixture mapping in admixed populations.
SparsePainter~\citep{yang2025sparsepainter} achieves scalability to biobank-size
cohorts through PBWT-based sparse haplotype matching.
\AP{} translates this established LAI architecture to the archaic reference
setting, replacing the modern reference panel with the Vindija or Altai archaic
genomes and reparameterizing the emission model for the coalescent depth of
archaic divergence.

\subsection*{Functional and adaptive significance of introgressed segments}

Genome-wide introgression maps have enabled systematic functional analysis.
These analyses depend on accurate, per-haplotype introgressed segment maps:
methods that mislocalize or miss short tracts may misattribute the source
of functional enrichment.
Genome-wide analyses have revealed Neanderthal allele enrichment in keratin and
immune pathways and depletion near promoters~\citep{sankararaman2016combined},
population-specific metabolic and neurological pathway enrichment in East Asian
and South Asian populations~\citep{vernot2016excavating}, phenotypic associations
in the UK Biobank~\citep{dannemann2017contribution}, and altered transcript levels
in immune cells for Neanderthal haplotypes at immune
gene loci~\citep{mccoy2017impacts}.
Enard \& Petrov~\citeyear{enard2018human} demonstrated global excess of archaic
introgression near antiviral immune genes, arguing for adaptive retention of
archaic viral defense alleles following exposure to Eurasian pathogens.
The functional annotation results of the present work --- the unsupervised
recovery of the interferon receptor gene cluster in CEU haplotypes --- provide
independent corroboration of this adaptive retention hypothesis and demonstrate
that \AP{}'s match-length signal is specific enough to locate these loci
without annotation input.

\subsection*{Ancient DNA and reference genome quality}

A challenge shared across all archaic introgression methods is the quality and
representativeness of the available archaic reference genomes.
The high-coverage Vindija 33.19 Neanderthal
genome~\citep{prufer2014complete, prufer2017high} is the best available
Neanderthal reference, but represents a single individual from one
geographic locality and time point.
The Altai Denisovan~\citep{meyer2012high} similarly represents a single
individual.
Ref genome bias, post-mortem DNA damage, and mapping artifacts can introduce
systematic errors in any method that compares modern and archaic sequences.
\AP{} addresses the reference-divergence problem via the positive-only emission
mode, which allows for uncertainty in whether a specific site in the reference
reflects the true introgressing allele.
This design choice is consistent with approaches in ancient DNA analysis that
use asymmetric likelihood functions to handle deamination-induced
transitions~\citep{jonsson2013mapdamage}.
